\chapter{Causal Relationships}

\section{Association versus causation}

\begin{definition}
Variates $x$ and $y$ are \emph{associated} if the distribution of $y$ differs for
different values of $x$ (e.g.\ a non-zero correlation, a non-zero regression slope, or
dependence in a two-way table). We say $x$ \emph{causes} $y$ if, other things being held
fixed, changing the value of $x$ for a unit changes the (distribution of the) value of $y$
for that unit.
\end{definition}

Association is about what we \emph{observe}; causation is about what would happen if we
\emph{intervened}. All the tests of Chapters 5--7 detect association. Whether we may
interpret it causally depends on how the data were collected.

\begin{example}[Associations that are not causal]
\begin{itemize}
  \item Cities with more firefighters have more fire damage (both are driven by city
        size and number of fires).
  \item Ice cream sales and drownings are positively correlated (both increase in hot
        weather).
  \item Students who use the tutoring centre have lower grades (weaker students are more
        likely to seek help: \emph{reverse causation} / selection).
\end{itemize}
\end{example}

\section{Confounding}

\begin{definition}[Confounding variate]
A variate $z$ is a \emph{confounder} for the relationship between $x$ and $y$ if $z$ is
associated with $x$ and also affects $y$. Then an observed association between $x$ and
$y$ may be partly or wholly due to $z$, and the effects of $x$ and $z$ cannot be separated
from the data.
\end{definition}

\begin{center}
\begin{tikzpicture}[node distance=2.2cm, v/.style={draw, circle, minimum size=0.9cm}]
  \node[v] (z) {$z$};
  \node[v, below left=1.2cm and 1.4cm of z] (x) {$x$};
  \node[v, below right=1.2cm and 1.4cm of z] (y) {$y$};
  \draw[-{Stealth}] (z) -- (x);
  \draw[-{Stealth}] (z) -- (y);
  \draw[-{Stealth}, dashed] (x) -- node[below] {\small causal effect?} (y);
\end{tikzpicture}
\end{center}

Strategies for dealing with confounding:
\begin{itemize}
  \item \textbf{Randomization} (in experiments): breaks any link between $x$ and every
        possible confounder, known or unknown.
  \item \textbf{Stratification/matching}: compare $x$-groups within levels of a known
        confounder $z$ (e.g.\ within age groups); pairing (Chapter 6) is an example.
  \item \textbf{Including $z$ in the model} (multiple regression, beyond Stat 231).
\end{itemize}
Only randomization handles \emph{unknown} confounders.

\section{Experimental versus observational studies}

\begin{definition}
In an \emph{experimental study} the investigator controls (assigns) the value of the
explanatory variate $x$ for each unit. In an \emph{observational study} the values of $x$
are not controlled; they are just recorded as they occur.
\end{definition}

\begin{important}[Principles of a good experiment]
\begin{itemize}
  \item \textbf{Control/comparison}: compare the treatment with a control group (e.g.\
        placebo), keeping other conditions the same.
  \item \textbf{Randomization}: assign treatments to units at random, so that the groups
        are alike on average in all other respects; then a difference in $y$ can be
        attributed to $x$ (or chance, which the p-value quantifies).
  \item \textbf{Replication}: use enough units to make the sample error small.
  \item \textbf{Blinding}: if possible, neither units (single-blind) nor evaluators
        (double-blind) know which treatment was received, avoiding placebo effects and
        measurement bias.
\end{itemize}
\end{important}

For many important questions (smoking and cancer, pollution and asthma) experiments are
unethical or impossible, so we must rely on observational studies: \emph{cohort} studies
(follow exposed and unexposed groups forward in time) and \emph{case-control} studies
(compare past exposure of people with and without the disease). Such studies can be very
informative but are always open to confounding.

\section{Simpson's paradox}

\begin{definition}
\emph{Simpson's paradox} occurs when an association between two variates in the
combined data has the opposite direction from the association within every subgroup
defined by a third variate. It is an extreme form of confounding.
\end{definition}

\begin{example}[Kidney stone treatments]
Success rates of two treatments in an observational study:
\begin{center}
\begin{tabular}{lccc}
\toprule
 & small stones & large stones & all patients\\
\midrule
Treatment A & $81/87 = 93\%$ & $192/263 = 73\%$ & $273/350 = 78\%$\\
Treatment B & $234/270 = 87\%$ & $55/80 = 69\%$ & $289/350 = 83\%$\\
\bottomrule
\end{tabular}
\end{center}
A is better for small stones and better for large stones, yet B looks better overall. The
explanation: stone size is a confounder. Large stones are harder to treat (lower success
for both treatments), and doctors gave A mostly to large-stone patients ($263$ of $350$)
and B mostly to small-stone patients ($270$ of $350$). The overall comparison mixes the
treatment effect with the severity effect. The stratified comparison (within stone size) is
the relevant one, and a randomized trial would avoid the problem.
\end{example}

\begin{remark}
Algebraically, the overall rate is a weighted average of the subgroup rates, with weights
equal to the fraction of each group's patients in each subgroup:
$0.78 = \frac{87}{350}(0.931) + \frac{263}{350}(0.730)$ and
$0.826 = \frac{270}{350}(0.867) + \frac{80}{350}(0.688)$. Different weights can reverse
the order.
\end{remark}

\section{Establishing causation from observational data: Hill's criteria}

When experiments are impossible, the epidemiologist Austin Bradford Hill (1965) proposed
considerations that strengthen the case for a causal interpretation of an association.
Stat 231 emphasizes:
\begin{enumerate}
  \item \textbf{Strength}: the association is large (e.g.\ heavy smokers have many times
        the lung cancer rate of non-smokers); a strong association is harder to explain
        away by confounding.
  \item \textbf{Consistency}: the association is found repeatedly in different studies,
        populations, designs and by different investigators.
  \item \textbf{Specificity}: the cause leads to a specific effect rather than many
        unrelated ones.
  \item \textbf{Temporality}: the cause precedes the effect.
  \item \textbf{Dose--response (biological gradient)}: more exposure leads to more
        effect.
  \item \textbf{Plausibility}: there is a credible mechanism (e.g.\ carcinogens in
        smoke).
  \item \textbf{Coherence}: the causal interpretation does not conflict with known facts
        about the disease.
  \item \textbf{Experimental evidence}: e.g.\ animal experiments, or the effect diminishes
        when exposure is removed (people who quit smoking have reduced risk).
  \item \textbf{Analogy}: similar causes are known to have similar effects.
\end{enumerate}
No single criterion is decisive (except temporality, which is necessary), but together they
built the overwhelming case that smoking causes lung cancer without a randomized trial.

\begin{important}[Common pitfalls]
\begin{itemize}
  \item A small p-value says an association is unlikely to be due to chance alone; it
        says nothing about whether it is causal.
  \item Random \emph{sampling} lets you generalize to a population; random
        \emph{assignment} lets you conclude causation. Neither implies the other.
  \item Always ask ``what else differs between the groups?''
\end{itemize}
\end{important}

\section*{Summary}
\begin{itemize}
  \item Association $\ne$ causation; confounders, reverse causation and selection can
        produce associations.
  \item Randomized controlled experiments are the gold standard for causal conclusions
        because randomization balances all confounders on average.
  \item Simpson's paradox: aggregating over a confounder can reverse an association.
  \item For observational data, Hill's criteria (strength, consistency, temporality,
        dose--response, plausibility, \ldots) support but never prove causation.
\end{itemize}
